\chapter{Geometry Governs Computational Capacity}
\label{chapter:act-III}


The previous chapter measured what a driven network's activity looks like, and
left two things unfinished. The first is that how many directions a readout can
use turned out to be a property of a substrate \emph{at an operating point}
rather than of the substrate alone (Section~\ref{sec:act2-floor}), so a
network's capability is a curve and has to be read across its range. The second
is that the common mode was offered as what those differences are made of, and
never shown to be: every substrate compared there carries non-negative weights,
so every one of them carries a common mode, and no comparison among them can
isolate it. This
chapter reads the curve, and then takes the mode away.

Two demands are put to these networks. Memory asks a network to reconstruct an
input it has already been given, which requires that traces of the past remain
separable in the present state. Prediction asks it to sustain a signal with
nothing arriving to correct it, the network driving itself from its own output.
The memory-capacity task puts the first demand directly, scoring the fraction of
a noise input a readout recovers at each of a range of delays and summing across
them, so that capacity is read in units of delays retained. Closed-loop
prediction of the chaotic Lorenz system puts the second, and
Chapter~\ref{chapter:methods} gives the details of both protocols.

Memory separates these substrates, but not where it was looked for. At its own
best operating point each one is using very nearly every direction available to
it, so peak capacity cannot be resolved at this size, and what can be compared
there amounts to parity. What separates them is what survives being driven away
from that point, and there the ordering reverses: the substrate whose peak sits
lowest of the four retains by far the most. The advantage is a width rather than
a height.

Removing the common mode says what that width is made of. Making a fraction of
the weights negative dissolves the guarantee that produces the mode, leaving the
topology and the weight magnitudes as they were, and the memory ladder closes
with it. It does not close in the direction the word removal suggests. No
substrate is harmed by the intervention; supercritical memory rises for all four,
and the advantage disappears because the nulls catch up from far below. What the
connectome's spectral gap supplies is therefore a rescue from a handicap that
non-negativity imposes on everyone, and not a capacity of its own.

The same intervention read on prediction does the opposite. Negative weights
bring a second dynamical branch within reach, and a network falls onto one branch
or the other rather than moving between them: a single bit recording which one
accounts for almost everything the full trajectory geometry accounts for. An
advantage opens where none existed, and across a band of sign fractions the
connectome is the only substrate still predicting usefully. One structural
property is therefore read by the two demands with opposite sign: the sign
composition that buys one costs the other, on every substrate tested.

One thing leaves this chapter unresolved, and it sits inside the account rather
than beside it. On the non-negative weights the instrument actually produces,
prediction is lost across the operating range with no geometric event of any
kind, so the mechanism this chapter assembles is built on a quantity that
registers nothing in exactly the regime the substrate occupies. What governs
that loss is named as an open problem and left that way.



\section{Memory beyond the peak}
\label{sec:act3-memory-sigma}
\bigskip

The primary variable of this chapter is the operating point: how hard the
recurrent dynamics are driven, set by rescaling each connectivity matrix to a
nominal spectral radius $\sigma$. Section~\ref{sec:act1-handover} sets out the
two axes that rescaling makes available and why neither of them is neutral. What
is worth restating here is the consequence for this chapter's claim in
particular: the matched-bulk axis leaves the Perron root unmatched, and so hands
the connectome more of the very quantity proposed as the cause. Every result
below names the axis it was measured on.

Memory is read here through the ridge effective rank of the driven states,
$d_{\mathrm{eff}}$, which counts how many directions of their un-centred Gram
matrix stand clear of the readout's regularisation floor and are therefore
available to a linear map. Chapter~\ref{chapter:act-II} established it as the
quantity that tracks measured memory capacity across substrates, where the
variance-weighted count does not, and it is used in place of the task score
wherever the argument turns on how many directions exist rather than on what a
particular readout recovers from them. Because a network of $N$ units has $N$ directions to offer,
$d_{\mathrm{eff}}$ is bounded above by $N$.

Matched on the bulk radius, all four substrates rise together, all four come
within a few percent of that bound, and they separate only as the dynamics are
driven further into the supercritical regime (Figure~\ref{fig:F7}). The connectome's peak is the lowest
of the four, at
$432.4$ of a possible $448$. % TIER0 1.2
At the top of the range every substrate and every seed covers, it holds
$47\%$ % TIER0 1.2
of that peak against the nulls'
$28$, $22$ and $11\%$. % TIER0 1.2
The ordering at the peak and the ordering deep in the supercritical regime are
opposite, and the crossing between them, rather than either ordering alone, is
the result. That crossing is not an artefact of the axis that produces it, since
the same retention ordering holds on the nominal axis, % TIER0 1.2
but neither is it a global advantage: on that axis it runs the other way below
criticality, and it is absent near the peak on both. Where the bulk is held
fixed, the defensible summary is parity below criticality and advantage above it.

\begin{figure}[htbp]
    \centering
    \includegraphics[width=\textwidth]{F7.png}
    \caption[Matched on the bulk radius, the connectome peaks lowest and retains
    most]{\textbf{Matched on the bulk radius, the connectome peaks lowest and
    retains most: the memory result is a crossing, not a capacity difference.}
    Memory-capacity task, $N = 448$, non-negative weights, ten seeds, ridge
    $\alpha = 10^{-6}$. \textbf{(a)} Ridge effective rank against
    $\sigma \cdot \mathrm{bulk95}$, the axis that holds the bulk radius fixed.
    $\mathrm{bulk95}$ is a per-seed quantity for the three resampling nulls, so each
    seed is reindexed onto a common grid before the median is taken; the dashed rule
    at $2.599$ marks the top of the range every substrate and every seed covers. All
    four rise together to the $d_{\mathrm{eff}} = N = 448$ ceiling and separate only
    on the way down. \textbf{(b)} At the peak they are not separable: each lands
    within $3.5\%$ of the ceiling, the connectome at $0.965$ against the nulls'
    $0.993$ to $0.997$. \textbf{(c)} Retention is separable and is the result: at the
    rule in (a) the connectome holds $47\%$ of its own peak against $28$, $22$ and
    $11\%$. A ceiling can clip curves but cannot manufacture a crossing.}
    \label{fig:F7}
\end{figure}

That the peak cannot carry a claim is a limit of the measurement rather than a
retreat from it. Each substrate at its own optimum is using very nearly every
direction available to it, within a few percent of the ceiling, % TIER0 1.2
and the larger network reported below does not escape that ceiling either.
What is resolvable is what happens as the networks are driven past that point,
and a difference there is robust to finite size in a way a difference between the
optima is not.

The peaks can still be compared, and the comparison runs against the chapter's
direction. Paired within seed, each substrate at its own optimum and each pair
sharing input weights and input series, the connectome's peak memory capacity
sits
$2$ to $6\%$ % TIER0 3.4
below the nulls'. Against degree-matching the difference is not consistently
signed, and against the other two nulls it is small and consistent.
Appendix~\ref{app:peak-parity} gives the comparison in full, across five ridge
strengths, and shows there that the ordering reported by Aceituno et
al.~\cite{aceituno2020tailoring} reproduces exactly on these substrates at the
smallest ridge. The defensible statement is parity at the peak rather than a
deficit: the crossing above is not a weak network falling from a lower start.

Whether the result is an artefact of the parcellation is settled by repeating the
memory sweep at
$N = 1000$, % TIER0 2.4
read on measured memory capacity rather than on ridge effective rank
(Figure~\ref{fig:F9}). The supercritical margin holds, and what counts as
supercritical decides how it should be read. Applying the connectome's own
critical scale to every variant, the connectome-to-Erdős--Rényi margin is
$4.40$ at the smaller size and $4.42$ at the larger; % TIER0 2.4
taking each variant above its own critical scale instead gives
$3.56$ and $3.85$. % TIER0 2.4
The first filter leads because it evaluates every substrate over an identical
range of operating points, and it is also the one that flatters the
connectome, since it samples each null further above its own critical point than
the null's own scale would. What survives both is that the margin is not an
artefact of the smaller network; what does not is calling it scale-invariant
without naming the filter. Appendix~\ref{app:scale} gives the protocol, the
control separating the two changes in play, and the one pre-registered test this
run was designed to settle and did not.

\begin{figure}[ht]
    \centering
    \includegraphics[width=\textwidth]{F9.png}
    \caption[The supercritical memory margin holds across a 2.2-fold change in
    $N$, under both definitions of ``supercritical'']{\textbf{The supercritical
    memory margin holds across network scale, under both definitions of
    ``supercritical''.} Memory-capacity task, non-negative weights, four
    substrates and ten seeds each; the connectome and Erdős--Rényi are drawn,
    since the panel is about scale rather than about the ladder. Run length was
    scaled from $3000$ to $6000$ to hold retained timesteps per unit near $5.5$,
    and the ridge reparameterised to scale with the Gram matrix rather than held
    at an absolute value; Appendix~\ref{app:scale} separates the two changes.
    \textbf{(a)} Median memory capacity against nominal $\sigma$ at both sizes;
    the dotted rules are the connectome's critical scale, $3.078$ at $N = 448$
    and $3.985$ at $N = 1000$. \textbf{(b)} With that threshold applied to every
    variant, the connectome to Erdős--Rényi margin is $4.40$ and $4.42$, on
    levels $12.32$ against $2.80$ and $13.93$ against $3.15$. \textbf{(c)} With
    each variant taken above its own critical scale, $3.56$ and $3.85$. Both are
    drawn because the answer depends on the filter: the margin holds either way,
    but only (b) is flat. The larger network does not escape the
    $d_{\mathrm{eff}}$ ceiling.}
    \label{fig:F9}
\end{figure}

What accounts for the retention difference has already been measured.
Chapter~\ref{chapter:act-II} found the connectome the least dominated of the four
by its common mode at every supercritical operating point, despite carrying by
far the largest Perron root, and found its Gram spectrum standing correspondingly
further clear of the readout's floor. A large gap lets the leading mode be driven
hard without the bulk of the spectrum following it, so the activity that is not
the common mode survives, and what survives clear of the floor is what the
readout can use. The retention difference in Figure~\ref{fig:F7} is that gap,
read out. That is not a larger capacity, and nothing above claims one; it is a
capacity that survives being driven away from the conditions it was measured
under.




\section{Prediction and its two failures}
\label{sec:act3-prediction-sigma}
\bigskip

A network scored on memory is handed the true signal at every step. A network
asked to generate is released after a short initialisation and runs on its own
output, so its mistakes become its own next input and accumulate, and what is
scored is how long it stays with the true signal before leaving it
(Chapter~\ref{chapter:methods}). At the operating point where each substrate is
at its best, near criticality at $\sigma = 2$, the four are not
distinguishable: the connectome's
margins over the three nulls are
$+0.28$, $-0.01$ and $+0.44$ % TIER0 2.6
Lyapunov times, none of them significant. That is the same shape as the memory
result, whose four peaks also sat within a few percent of one another. On both
demands the substrates separate away from their best operating points rather than
at them.

They separate in two different ways, and only one of them favours the connectome.
Driven well into the supercritical regime, above
$\sigma \approx 7.6$ to $8.0$, % TIER0 2.3
five of Erdős--Rényi's ten seeds settle onto a pair of states and alternate
between them, flipping from one to the other at every step and going nowhere
else. What they emit is no longer a path winding around the Lorenz attractor; it
is two points (Figure~\ref{fig:F21}a). Across the whole swept range, far above
its own critical scale, this happens to none of the connectome's ten seeds. The substrates are not being
separated here by how long they predict for, but by whether what they produce is
a trajectory at all. Five against none is significant
($p = 0.033$, % TIER0 2.3
Fisher exact), though on ten seeds per substrate the claim rests on the size and
direction of the difference rather than on the test.

\begin{figure}[ht]
    \centering
    \includegraphics[width=\textwidth]{F21.png}
    \caption[Prediction falls across the operating-point sweep while trajectory
    curvature does not move]{\textbf{Prediction falls across the
    operating-point sweep while the geometry does not move.} Lorenz task,
    $N = 448$, non-negative weights, ten seeds per substrate, against nominal
    $\sigma$. Points are individual seeds and lines are seed medians.
    \textbf{(a)} Trajectory curvature, the mean angle between successive segments
    of the driven trajectory, drawn on the full $0$ to $\pi$ range so that
    flatness is not an artefact of a compressed axis. Seeds occupy two separated
    groups, one near $0.26$~rad and one just below $\pi$, and none lies between
    them. The upper group is the alternating state of the text: a two-state cycle
    reverses direction at every step, which this measure reads as an angle near
    $\pi$. Medians mislead here, and Erdős--Rényi's $1.70$~rad at
    $\sigma = 11.2$ is a value that no seed takes. \textbf{(b)} Valid prediction
    time in Lyapunov times, log axis, falling for every substrate across the same
    range; seeds scoring exactly zero are stacked on the strip below the break,
    where height reads as a count. Curvature is measured on the driven trajectory
    and prediction time on the free run. The graded loss in (b) has no
    counterpart in (a).}
    \label{fig:F21}
\end{figure}

What the connectome brings to that comparison is the same spectral gap the memory
result turned on, read from its other end. Alternation between two states is what
a network settles into when the most negative part of its spectrum grows strong
enough to dominate its own dynamics. A large gap is exactly what holds that part
far below the leading eigenvalue, so the substrate with the largest gap is the
one furthest from reaching the alternating state, and it stays furthest from it
however hard it is driven. As stated, that is an account that fits rather than
one that has been tested: every substrate compared so far has the same sign
structure, so no comparison among them can isolate it. The next section supplies
the intervention that can, and Section~\ref{sec:act3-unifying} returns to the
account itself.

The second way of failing does not favour the connectome. Counted bluntly, by how
many seeds cease to predict at all somewhere high in the sweep, the ordering
reverses: the connectome is the worst of the four, and the degree-matched null
never fails at any operating point. The two failures are separate events rather
than degrees of one, and can be shown to be: the connectome has no seed in the
alternating state anywhere, so none of the seeds on which it ceases to predict
can be a seed that collapsed. Appendix~\ref{app:second-failure} gives the
counts. What holds is the narrower statement, that the connectome resists one
specific way of failing to generate.

Neither failure shows up in the shape of the activity. The connectome's
curvature, the mean angle between successive segments of its trajectory, holds at
$0.26$~rad % TIER0 3.11
from just above criticality to the top of the sweep, while over the same range
its prediction falls roughly tenfold (Figure~\ref{fig:F21}). The measure that
carried the memory half of this chapter registers nothing here: capacity is lost
with the geometry intact. What governs that loss is left open in
Section~\ref{sec:act3-unifying}, which states the three explanations available
and eliminates each.

% Every substrate to this point has non-negative weights, so every one of them
% carries a common mode, and no contrast between them can show that the common mode
% is what any of these differences are made of. The next section stops comparing
% and starts changing the weights.





\section{Inhibition and the common mode}
\label{sec:act3-lesion}
\bigskip

Neither of the last two sections can settle its own mechanism, because every
substrate compared in them is non-negative and therefore carries a common mode: a
comparison among four networks that all share a feature cannot show that the
feature is what separates them. The account offered was that a large spectral gap
lets the leading mode be driven hard without the bulk of the spectrum following
it. If that is right, taking the common mode away should take the difference with
it. A fraction $f$ of the weights is therefore made negative, leaving the
magnitudes untouched and the topology intact, so that what changes between one
network and the next is the sign pattern and nothing else. As $f$ approaches one
half the excitatory and inhibitory contributions to each unit's input come into
balance and the all-positive leading direction has nothing left to be built from.
This is not a perturbation away from biology. Macroscale diffusion tractography
cannot represent the sign of a connection, so a non-negative matrix is what the
instrument produces rather than what cortex is, and raising $f$ asks what a
network carrying inhibition would do.\footnote{Which edges are inverted is
controlled rather than left to chance: flips are stratified by endpoint
degree, so that the inverted set has the same hub-to-periphery composition as
the graph at every $f$ and the amount of inhibition does not vary with where it
sits. Chapter~\ref{chapter:methods} gives the procedure.}

The intervention does what it was designed to do, and does it to every substrate
alike. Held at the supercritical operating point
$\sigma = 6$, % TIER0 3.7
where Section~\ref{sec:act3-memory-sigma} found the substrates furthest apart,
the common-mode amplitude falls for all four as $f$ rises, and by the time the
signs are balanced it has fallen by two orders of magnitude on every one of them
(Figure~\ref{fig:F11}a). The four curves converge as they fall, because the
connectome starts nearest the floor and has least distance to travel: the nulls
lose a feature that dominated their activity, and the connectome loses one it was
already largely free of. What the sign balance removes is therefore not something
the connectome had and the nulls lacked. It is the handicap the nulls carried and
the connectome did not.

\begin{figure}[htbp]
    \centering
    \includegraphics[width=\textwidth]{F11.png}
    \caption[Non-negativity makes memory worse for every substrate, and the
    connectome's spectral gap makes it least worse]{\textbf{The connectome does
    not make memory better; non-negativity makes it worse for every substrate,
    and the connectome's spectral gap makes it least worse.} Memory-capacity
    task, $N = 448$, ten seeds, held at the supercritical operating point
    $\sigma = 6$ throughout, against the fraction $f$ of weights made negative.
    \textbf{(a)} The common-mode amplitude, the grand mean of the state matrix
    over time and units, absolute-valued per cell before the seed median. With
    non-negative weights the nulls sit at $0.949$ to $0.989$, very nearly every
    unit pinned at one saturated value, while the connectome sits at $0.759$
    despite carrying much the largest Perron root; balancing the signs removes
    the mode entirely for every substrate. \textbf{(b)} The median absolute
    connectome-minus-Erdős--Rényi gap in memory capacity under each matching
    axis, falling from $6.42$ to $4.75$ with non-negative weights and to near
    $0.5$ under either axis by $f = 0.2$.}
    \label{fig:F11}
\end{figure}

The memory ladder goes with it, and the way it goes is the finding. The wedge
between the connectome and the nulls narrows as $f$ rises and has closed by
$f \approx 0.15$ to $0.20$. % TIER0 2.6
No substrate is harmed in the process. Supercritical memory capacity rises with
$f$ for all four, the connectome's from
$11.43$ to $14.18$ % TIER0 2.6
and Erdős--Rényi's from
$2.42$ to $13.11$ % TIER0 2.6
across the swept range at the same operating point. The advantage disappears
because the nulls catch up, and they catch up from far below because they had
much further to come. That is the prediction the common-mode account makes and
the alternatives do not: on a reading in which the connectome's weight assignment
confers some intrinsic computational property, balancing the signs should have
degraded the substrate that had the property, and instead it lifted the ones that
lacked it. The supercritical advantage is a rescue from a handicap that
non-negativity imposes on everyone, and not a capacity the connectome supplies.

The intervention also bounds how much of the ladder the bulk radius accounts for,
a question Section~\ref{sec:act3-memory-sigma} raised and could not settle.
At $f = 0$, matching the substrates on the bulk radius rather than on the Perron
root absorbs about
a quarter % A3M 2.6
of the connectome-to-Erdős--Rényi gap; once the signs are balanced it absorbs
almost none of it (Figure~\ref{fig:F11}b). It is a partial
controller, and it explains something only where there is a common mode to
explain.




\section{A second dynamical regime}
\label{sec:act3-branch}
\bigskip

Applied to prediction, the same intervention does something else entirely, and
seeing why takes one fact about what a negative weight does dynamically. Each
unit passes the weighted sum of its recurrent input, plus the drive, through a
$\tanh$. A map of that shape settles to a fixed point when its effective gain is
sufficiently positive and to a period-two orbit when the gain is sufficiently
negative, with nothing stable between the two. Non-negative weights can only
reach the first; making a fraction of the weights negative brings the second
within reach. The intervention therefore does not turn a dial. It makes a second
dynamical regime available, and a network either enters it or does not.

The activity behaves as the map says it should, and the two-group structure that
appeared at $f = 0$ in Section~\ref{sec:act3-prediction-sigma} now appears across
the whole sweep. Trajectory curvature occupies two narrow spikes with almost
nothing between them:
$98\%$ % TIER0 3.9
of cells sit either near
$0.25$~rad, % TIER0 3.9
where successive steps continue in roughly the same direction, or at
$\pi$, % TIER0 3.9
where they reverse, which is what a period-two orbit looks like measured this way
(Figure~\ref{fig:F12}). What matters for capacity is which spike a network is
in, not where on the curvature axis it sits: a single binary variable recording
the spike explains as much of the variation in prediction as the full continuous
measure does. This also withdraws a graded account that an earlier analysis of
these data supported, under which straighter trajectories predicted better;
pooled across all
cells the correlation is strong, but it is an artefact of mixing the two clusters
and it reverses in sign within either one.
Appendix~\ref{app:regime-gating} gives the regression and the reversal. Capacity
is not gated by how curved the manifold is. It is gated by which of two dynamical
regimes the manifold is in.

\begin{figure}[htbp]
    \centering
    \includegraphics[width=\textwidth]{F12.png}
    \caption[Capacity is gated by dynamical regime, not graded by trajectory
    curvature]{\textbf{Capacity is gated by dynamical regime, not graded by
    trajectory curvature.} Lorenz task, $N = 448$, all four substrates over the
    full sweep in operating point and sign fraction. \textbf{(a)} Distribution of
    trajectory curvature. Two spikes, one near $0.25$~rad and one at $\pi$ where
    successive steps reverse, hold $98\%$ of cells between them; under $0.6\%$
    fall in the valley, whose midpoint at $1.0$~rad is used as a separator.
    \textbf{(b)} Prediction against curvature, coloured by cluster: a binary
    collapsed-or-not variable explains $R^2 = 0.364$ of the variance against
    continuous curvature's $0.370$. \textbf{(c)} Within the smooth cluster only,
    where the residual correlation is $+0.145$, opposite in sign to the pooled
    $-0.78$. Appendix~\ref{app:regime-gating} gives both comparisons.}
    \label{fig:F12}
\end{figure}

Which regime a network falls into, and at what sign fraction, is where weight
assignment reappears, and it is where a generative advantage opens that did not
exist before. At the near-critical operating point $\sigma = 2$ the connectome is
level with all three nulls on non-negative weights, as
Section~\ref{sec:act3-prediction-sigma} reported. A margin opens as the sign
fraction rises (Figure~\ref{fig:F13}), and at
$f = 0.25$ % TIER0 2.6
the connectome's paired advantage in valid prediction time is
$+1.71$, $+1.62$ and $+2.20$ % TIER0 2.6
Lyapunov times over the weight-permuted, degree-matched and Erdős--Rényi nulls,
the first sign fraction at which it clears all three. Across the range from
$f \approx 0.20$ to $0.25$ % TIER0 2.6
the connectome is the only substrate still predicting usefully at all, at
$1.3$ to $2.8$ Lyapunov times % TIER0 2.6
against the nulls'
$0.1$ to $0.9$. % TIER0 2.6
What licenses reading this as a property of assignment rather than of topology is
the control it clears: the weight-permuted null holds the connectome's topology
and the connectome's own weights, and changes only which weight sits on which
edge. The margin is not monotone in $f$, and it narrows considerably if what
counts as predicting is raised from not having failed outright to predicting
usefully. It is real and it is located; it is not a uniform superiority across
the sweep.

\begin{figure}[htbp]
    \centering
    \includegraphics[width=\textwidth]{F13.png}
    \caption[A predictive advantage opens as signs are balanced, and clears the
    assignment control]{\textbf{A generative advantage opens as signs are
    balanced, and clears the assignment control.} Lorenz task, $N = 448$, ten
    seeds, at the near-critical operating point $\sigma = 2$ throughout, against
    the sign fraction $f$. \textbf{(a)} Median valid prediction time per
    substrate. \textbf{(b)} The connectome's advantage over each null, paired
    within seed so that every comparison shares its input weights and input
    series. None of the three margins is significant on non-negative weights; all
    three are cleared first at $f = 0.25$, and the margin is smaller at
    $f = 0.50$ than at its best, so it is not monotone in $f$. \textbf{(c)} At
    $f = 0$ and far above criticality, seeds collapsing to a period-two orbit:
    five of ten for Erdős--Rényi and none of ten for the connectome
    (Section~\ref{sec:act3-prediction-sigma}). This panel covers only the two
    substrates the boundary capture contains.}
    \label{fig:F13}
\end{figure}







% \newpage

\section{One property, two demands}
\label{sec:act3-unifying}
\bigskip

The two demands read one property with opposite sign, and it is the property
Chapter~\ref{chapter:act-I} measured. The connectome's spectral gap is the
largest on the ladder: its Perron root stands
$3.08$ % TIER0 3.1
times clear of its bulk edge against Erdős--Rényi's
$1.81$ % TIER0 3.1
(Section~\ref{sec:act1-spectrum}). Expressed the other way up, everything
beneath the Perron root sits proportionally nearer the origin, and because these
matrices are symmetric at $f = 0$ that includes the negative end: the most
negative eigenvalue reaches
$0.33$ % TIER0 3.9 [VERIFY quantity]
of the leading one for the connectome against Erdős--Rényi's
$0.55$. % TIER0 3.9 [VERIFY quantity]
So stated, the gap is a distance to the second branch of the map of
Section~\ref{sec:act3-branch}. A network settling toward the fixed-point branch
produces smooth trajectories, which is what a closed loop needs in order to keep
generating; the same settling synchronises the units, so they stop doing
different things and the readout dimensions a linear map could have used are
wasted. What helps prediction is what costs memory.

Put both demands on one grid and the reversal is visible directly
(Figure~\ref{fig:F22}). Memory improves as
the sign fraction rises, and the improvement is an ordering of the supercritical
decay: steepest on non-negative weights and nearly flat once a quarter of the
weights are negative, which is the rescue of Section~\ref{sec:act3-lesion} shown
as a family of curves rather than asserted. Prediction orders the other way, its
tallest peaks belonging to the lowest sign fractions and the sign-balanced curves
reaching zero well before the end of the sweep. What divides the two is not how
far back each reaches into the past but whether the network runs on its own
output, which is what makes the regime it settles into consequential in one case
and harmless in the other.

\begin{figure}[h!]
    \centering
    \includegraphics[width=\textwidth]{figures/fig_frontier_absolute.png}
    \caption[Memory and prediction respond to the sign fraction in opposite
    directions]{\textbf{Memory and prediction respond to the sign fraction in
    opposite directions, on every substrate.} $N = 448$, four substrates and ten
    seeds each, plotted against nominal $\sigma$; curves are seed medians, coloured
    by the fraction $f$ of weights made negative, from $f = 0$ (non-negative, dark)
    through $f = 0.5$ (sign-balanced, light). Columns are the four-substrate ladder.
    Dotted rules mark each substrate's own critical scale on its non-negative matrix,
    and shading marks the range of $\sigma$ over which most seeds sit at the metric's
    floor. \textbf{(a)--(d)} Memory capacity: the total input variance a linear
    readout recovers across all delays, higher being better, peaking between $14.3$
    and $15.5$. \textbf{(e)--(h)} Valid prediction time: how long a released network
    stays with the true Lorenz trajectory, in Lyapunov times, higher being better.
    Peaks reach $4.4$ for the connectome and degree-matching and $3.7$ for the other
    two; the sign-balanced curves peak below $1.8$ and reach zero from
    $\sigma \approx 4$ onward. Vertical scales are shared within each row and are not
    comparable between rows.}
    \label{fig:F22}
\end{figure}

The memory advantage of Section~\ref{sec:act3-memory-sigma} and the predictive
advantage of Section~\ref{sec:act3-branch} are therefore not two findings that
happen to point the same way. They are one structural property measured by two
demands that want opposite things from it, and what in that account is measured
should be separated from what is inferred. The measurements are the two-state
distribution of trajectory curvature, the fact that a binary regime variable
explains nearly as much of prediction as the full geometry does, the ordering of
common-mode amplitude across substrates, the operating points at which
free-running trajectories collapse, and the opposed responses of the two demands
to the sign fraction. The inference is that a single account, built on where a
substrate's leading eigenvalue sits relative to the rest of its spectrum,
generates all of it. That account is consistent with everything measured here, and it is not a
derivation: nothing here shows that the two readouts must follow from the one
cause, only that they do wherever this work has looked.

The account also has a boundary, and it runs through the middle of this chapter
rather than around its edge. Capacity gated by dynamical regime is a statement
about networks carrying negative weights, and not about the substrate the
instrument actually produces. A non-negative matrix has no dominant negative
eigenvalue to flip, so the second branch is unreachable and there is no switch to
be gated by, and the quantity the mechanism is built on registers nothing: this
is the flat curvature Section~\ref{sec:act3-prediction-sigma} measured against a
tenfold loss of prediction. Three explanations were available and each fails.
Not geometry, which does not move. Not the
effective spectral radius, which the transition does not track: where prediction
is being lost it stands two orders of magnitude below its own largest value and
on the descending branch of its own
curve. % TIER0 3.11 [VERIFY subject clause]
And not memory, since at $\sigma = 6$ the connectome holds
$11.43$ against Erdős--Rényi's $2.42$ % TIER0 3.11
in memory capacity and a \emph{lower} valid prediction time,
$0.81$ against $1.18$. % TIER0 3.11
Whatever governs prediction on non-negative weights, this work has not found it,
and no part of the account above should be read as covering it.

Read from the biological side, the two demands say something neither says
alone. Neither is served best by a network tuned to one point, and what the
connectome supplies on both is the width of the region in which it does not fail:
memory carried further from its optimum, and prediction sustained further into
inhibition than any null reaches. A structure organised that way is not one that
computes best. It is one that needs less precise tuning in order to compute at
all, which is the reading Chapter~\ref{chapter:discussion} returns to.
